% =====================================================================
%  Preamble for handout.tex
% =====================================================================
%  Sections in this file:
%    1. Fonts, language and page
%    2. Tables, math, graphics
%    3. Colours, figures and links
%    4. Code and example outputs
%    5. Bibliography, hyperref, glossary
%    6. Title page and small helpers
% =====================================================================

% The renders are PDF 1.7; without this, pdfTeX warns on every one of them.
\pdfminorversion=7

% ---------------------------------------------------------------------
%  1. Fonts, language and page
% ---------------------------------------------------------------------
\usepackage[T1]{fontenc}            % font encoding: T1
\usepackage[british]{babel}         % British hyphenation, dates and names
\usepackage{amsmath}                % before newtxmath, as its manual asks
\usepackage{newtxtext}              % Times for text
\usepackage{newtxmath}              % and for math (brings the AMS symbols)
\usepackage[varqu]{inconsolata}     % code font, with straight quotes
\usepackage{microtype}              % finer justification
\usepackage{csquotes}               % \enquote, and biblatex wants it with babel

% The text block is exactly as wide as the students' article
% (a4paper, 11pt: 360pt), so the example outputs (renders) are shown at
% their natural size and their text lines up with the handout's text.
\usepackage[a4paper, textwidth=360pt, top=3cm, bottom=3cm]{geometry}

\usepackage{parskip}                % no indent, a gap between paragraphs
% Long typewriter words (file names, commands) cannot be hyphenated;
% a little extra stretch lets TeX break the line before them instead.
\setlength{\emergencystretch}{1.5em}
\usepackage{enumitem}
\setlist[itemize]{itemsep=0ex, leftmargin=1em}

% ---------------------------------------------------------------------
%  2. Tables, math, graphics
% ---------------------------------------------------------------------
\usepackage{array}
\usepackage{booktabs}
\usepackage{tabularx}
\usepackage{multicol}               % cheat sheet
\usepackage{siunitx}

\usepackage{graphicx}
\usepackage[export]{adjustbox}      % frames around the example outputs

\usepackage{caption}
\captionsetup{%
    labelfont=bf,
    labelsep=quad,
}

% ---------------------------------------------------------------------
%  3. Colours, figures and links
% ---------------------------------------------------------------------
\usepackage{xcolor}
\definecolor{backcolour}{HTML}{f5f5f0}      % background of code
\definecolor{titlecolour}{HTML}{e4e4dc}     % title bar of a code box
\definecolor{commentcolour}{HTML}{6a9955}
\definecolor{keywordcolour}{HTML}{7f0055}
\definecolor{stringcolour}{HTML}{9e22bd}
\definecolor{numbercolour}{HTML}{808080}
\definecolor{linkcolour}{HTML}{1a4fbf}

% Colours of the figures shared with the slides (figures/*.tex).
\definecolor{figaccent}{HTML}{3d3dff}
\definecolor{figalt}{HTML}{64cf4e}
\definecolor{figfill}{HTML}{d1d1d1}

\usepackage{hologo}                 % \hologo{MiKTeX}
\usepackage{skipink}                % underlines that skip descenders

\usepackage{tikz}
\usetikzlibrary{%
    positioning,
    calc,
    arrows.meta,
    fit,
}
\tikzset{>=stealth}
\usepackage{pgfplots}
\pgfplotsset{compat=1.18}
\usepackage[edges]{forest}          % directory trees

% inline tikz node for later referencing
\newcommand{\tn}[2]{% from https://tex.stackexchange.com/a/402466/121799
    \ifmmode%
    \tikz[remember picture,baseline= (#1.base),inner sep=0pt]{%
        \node(#1){$#2$};
    }%
    \else%
    \tikz[remember picture,baseline= (#1.base),inner sep=0pt]{%
        \node(#1){#2};%
    }%
    \fi%
}

% ---------------------------------------------------------------------
%  4. Code and example outputs
% ---------------------------------------------------------------------
%  Code is set by listings, so the handout compiles without -shell-escape.
%  listings has no complete list of LaTeX commands, so the `ltx' language
%  below colours everything that starts with a backslash instead.
\usepackage{listings}
\lstdefinelanguage{ltx}{%
    alsoletter={\\@*},
    keywordsprefix={\\},
    morekeywords={\\begin},         % keywordsprefix needs one keyword to work
    morecomment=[l]{\%},
    % Escaped characters: not a command, and \% does not start a comment.
    literate=%
        {\\\%}{{\textbackslash\%}}2
        {\\\{}{{\textbackslash\{}}2
        {\\\}}{{\textbackslash\}}}2
        {\\\&}{{\textbackslash\&}}2
        {\\\#}{{\textbackslash\#}}2
        {\\\$}{{\textbackslash\$}}2
        {\\\_}{{\textbackslash\_}}2
        % listings cannot read UTF-8 by itself; the letters used in examples:
        {ä}{{\"a}}1 {ö}{{\"o}}1 {ü}{{\"u}}1
        {Ä}{{\"A}}1 {Ö}{{\"O}}1 {Ü}{{\"U}}1 {ß}{{\ss}}1,
}
% .bib files: the entry type (@book) is coloured.
\lstdefinelanguage{bib}{%
    alsoletter={@},
    keywordsprefix={@},
    morekeywords={@book},
    morecomment=[l]{\%},
}
\lstdefinestyle{code}{%
    language=ltx,
    basicstyle=\ttfamily\footnotesize,
    aboveskip=0pt,                  % the box has its own padding
    belowskip=0pt,
    columns=fullflexible,
    keepspaces,
    showstringspaces=false,
    tabsize=4,
    keywordstyle=\color{keywordcolour},
    commentstyle=\color{commentcolour},
    numbers=left,
    numberstyle=\ttfamily\scriptsize\color{numbercolour},
    numbersep=12pt,                 % line numbers sit in the margin
    breaklines,
    postbreak=\mbox{\textcolor{numbercolour}{$\hookrightarrow$}\space},
}
\lstset{style=code}
% Commands for a terminal (no colours, no line numbers):
%   \begin{codeblock}[style=shell] ... \end{codeblock}
% listings cannot read non-ASCII characters; quoted error messages that
% contain one write it as LaTeX between (* and *), e.g. (*$\leq$*).
\lstdefinestyle{shell}{%
    style=code,
    language={},
    keywordstyle={},
    commentstyle={},
    numbers=none,
    escapeinside={(*}{*)},
}

% Code boxes: light background, no frame.
% A box sticks out 8pt into both margins, exactly like the frame of a
% render (standalone's border), so code and output line up.
\usepackage{tcolorbox}
\tcbuselibrary{listings, skins, breakable, raster}
\newlength{\exoverhang}
\setlength{\exoverhang}{8pt}
\tcbset{%
    codebox/.style={%
        enhanced,
        listing only,
        colback=backcolour,
        frame hidden,
        boxrule=0pt,
        arc=1mm,
        boxsep=0pt,
        left=5pt, right=2pt, top=4pt, bottom=4pt,
        grow to left by=\exoverhang,
        grow to right by=\exoverhang,
        colbacktitle=titlecolour,
        coltitle=black,
        fonttitle=\small\bfseries,
        toptitle=2pt, bottomtitle=2pt,
        % No page break before a code box: it stays with the sentence that
        % introduces it, and with the code box or render before it.
        before={\par\nopagebreak\vskip\smallskipamount\nopagebreak},
        after={\par\nopagebreak\vskip\smallskipamount},
    },
}

% \codefile[<listings options>]{<file>}: a file as a code box.
% \codefiletitled[<listings options>]{<file>}{<title>}: the same, with a
% title bar (Preamble, Document, references.bib, ...).
% For a .bib file: \codefile[language=bib]{...}.
\NewTCBInputListing{\codefile}{O{} m}{%
    codebox,
    listing file={#2},
    listing options={style=code, #1},
}
\NewTCBInputListing{\codefiletitled}{O{} m m}{%
    codebox,
    title={#3},
    listing file={#2},
    listing options={style=code, #1},
}
% Code typed directly in the handout:
%   \begin{codeblock}[<listings options>] ... \end{codeblock}
% !O{}: the optional argument must follow \begin{codeblock} directly;
% otherwise the search for it would read the first line of code too early.
\DeclareTCBListing{codeblock}{!O{}}{%
    codebox,
    listing options={style=code, #1},
}
% The same with a title bar:
%   \begin{codeblocktitled}{<title>}[<listings options>] ... \end{codeblocktitled}
\DeclareTCBListing{codeblocktitled}{m !O{}}{%
    codebox,
    title={#1},
    listing options={style=code, #2},
}
% Long listings (the appendix) may break across pages.
\NewTCBInputListing{\codefilelong}{O{} m}{%
    codebox,
    breakable,
    listing file={#2},
    listing options={style=code, #1},
}

% \render[<includegraphics options>]{<file>}: a pre-compiled output.
% Renders are 360pt wide plus an 8pt border on each side, shown at their
% natural size: the text in them has the size the students see.
\newcommand{\render}[2][]{%
    \par\nopagebreak                      % stays with the code above
    \noindent\makebox[\linewidth][c]{%
        \includegraphics[cframe=black!35 0.4pt, #1]{#2}%
    }\par\medskip
}

% Where the slides are. Locally the handout sits next to them
% (../slides). Compiled from the repository root (Overleaf, see
% README.md), they are in slides/.
\IfFileExists{../slides/renders/snippet.tex}{%
    \newcommand{\slidesdir}{../slides}%
}{%
    \newcommand{\slidesdir}{slides}%
}

% Where the generated workshop template is (make in template/), found the
% same way. It is optional: without it, the appendix shows only the
% rendered pages.
\IfFileExists{../template/build/workshop-template/checkpoints/checkpoint-04.tex}{%
    \newcommand{\templatedir}{../template/build/workshop-template}%
}{%
    \newcommand{\templatedir}{template/build/workshop-template}%
}

% Shortcuts for the two places listings and renders come from:
%   \slidecode{02-itemize.tex}        slides/listings/02-itemize.tex
%   \sliderender{02-lists.pdf}        slides/renders/02-lists.pdf
%   \owncode{...}, \ownrender{...}    the handout's own listings/ and renders/
\newcommand{\slidecode}[2][]{\codefile[#1]{\slidesdir/listings/#2}}
\newcommand{\slidecodetitled}[3][]{\codefiletitled[#1]{\slidesdir/listings/#2}{#3}}
\newcommand{\sliderender}[2][]{\render[#1]{\slidesdir/renders/#2}}
\newcommand{\owncode}[2][]{\codefile[#1]{listings/#2}}
\newcommand{\owncodetitled}[3][]{\codefiletitled[#1]{listings/#2}{#3}}
\newcommand{\ownrender}[2][]{\render[#1]{renders/#2}}

% Inline code, coloured like the code boxes: \code|\section{Results}|
% Like all verbatim, it does not work inside the argument of another
% command (\section{...}, \caption{...}, \item[...]); use \texttt there.
\lstdefinestyle{inline}{%
    style=code,
    basicstyle=\ttfamily,
    numbers=none,
    breaklines=false,
}
\newcommand{\code}{\lstinline[style=inline]}

% A grey note under an example, like the grey lines on the slides.
\newcommand{\examplenote}[1]{{\small\color{black!60}#1\par}}

% ---------------------------------------------------------------------
%  5. Bibliography, hyperref, glossary
% ---------------------------------------------------------------------
\usepackage[backend=biber, style=numeric]{biblatex}
\addbibresource{handout.bib}

% hyperref last; the glossary comes after it, as glossaries must.
\usepackage{hyperref}
\hypersetup{%
    hidelinks,
    pdftitle={LaTeX Workshop: Handout},
}
\urlstyle{same}

% \autoref names, capitalised.
% babel sets the names when the language starts, so they are changed
% through its hook for british.
\addto\extrasbritish{%
    \def\chapterautorefname{Chapter}%
    \def\sectionautorefname{Section}%
    \def\subsectionautorefname{Section}%
    \def\subsubsectionautorefname{Section}%
    \def\appendixautorefname{Appendix}%
    \def\equationautorefname{Equation}%
}

% External links: blue and underlined, so they can be told apart from
% the internal references (black) on paper as well.
% The text is usually the address without https://, so a printed copy
% still says where the link goes.
\newcommand{\myhref}[2]{%
    \textcolor{linkcolour}{\skipul{\href{#1}{#2}}}%
}

% ------------------------------------------------------- acronyms/glossary --
% \makenoidxglossaries sorts the entries in LaTeX itself, so no extra
% program (makeglossaries, bib2gls) has to run.
\IfFileExists{tracklang.sty}{%
    \usepackage[abbreviations, toc]{glossaries-extra}
    \setabbreviationstyle[acronym]{long-short}
    \makenoidxglossaries
    \loadglsentries{appendix/glossary_entries.tex}
    \newcommand{\reportglossary}{\printnoidxglossary[nonumberlist, type=abbreviations, sort=word, title={Glossary}]}
}{%
    \newcommand{\reportglossary}{}
}

% ---------------------------------------------------------------------
%  6. Title page and small helpers
% ---------------------------------------------------------------------
\makeatletter
\renewcommand{\maketitle}{%
    \hypersetup{pageanchor=false}% the title page has no number of its own
    \begin{titlepage}
        \centering
        \vspace*{0.2\textheight}
        {\fontsize{50}{48}\selectfont\bfseries\@title\par}
        \vspace{1em}
        \ifx\@subtitle\@empty\else
            {\LARGE\@subtitle\par}
        \fi
        \vfill
        % {\large\@author\par}
        % \vspace{2em}
        {\large\@date\par}
    \end{titlepage}%
    \hypersetup{pageanchor=true}%
}
% Optional extra field, used like \title
\newcommand{\subtitle}[1]{\gdef\@subtitle{#1}}
\let\@subtitle\@empty
\makeatother

\newcommand{\bslash}{\char`\\}
\newcommand{\bropen}{\char`\{}
\newcommand{\brclose}{\char`\}}
